\title{DeepSeekMath: Pushing the Limits of Mathematical Reasoning in Open Language Models}

\begin{abstract}
Mathematical reasoning remains a formidable obstacle for language models, primarily because of its inherent complexity and structure. In this work, we present DeepSeekMath~7B, an extension of DeepSeek-Coder-Base-v1.5 7B that undergoes continued pre-training using 120B mathematics-focused tokens drawn from Common Crawl, alongside supplementary natural language and code corpora. Without incorporating external toolchains or ensemble methods, DeepSeekMath~7B attains a notable 51.7\% on the MATH benchmark at the competition level, nearing the effectiveness observed with Gemini-Ultra and GPT-4. Employing self-consistency over 64 generations from DeepSeekMath~7B, performance rises to 60.9\% on MATH. Two central innovations underpin the mathematical reasoning capabilities of DeepSeekMath: (1) a robust data selection pipeline that expertly leverages freely available web-based data, and (2) the development of Group Relative Policy Optimization (GRPO), an adaptation of Proximal Policy Optimization (PPO) designed to refine mathematical problem solving while simultaneously improving PPO’s memory efficiency.
\end{abstract}

\section{Introduction}

Large language models (LLMs) have transformed the landscape of mathematical reasoning in artificial intelligence, driving remarkable improvements in quantitative reasoning \citep{MATH} as well as geometric reasoning assessments \citep{trinh2024solving}. In addition, such models have increasingly served as valuable tools for aiding human efforts in tackling challenging mathematical questions \citep{tao}. Nonetheless, leading solutions like GPT-4 \citep{gpt4} and Gemini-Ultra \citep{gemini} remain closed-source, and publicly available open-source alternatives have yet to match their performance.

This paper presents DeepSeekMath, a specialized language model tailored for mathematical tasks, which demonstrates a notable advancement over existing open-source models and nearly reaches the proficiency of GPT-4 on academic mathematical evaluations. Central to this achievement is the construction of the DeepSeekMath Corpus, an extensive and high-quality pre-training resource consisting of 120B mathematical tokens. The dataset is sourced from Common Crawl (CC) by employing a classifier built on fastText \citep{joulin2016fasttext}. Initially, the classifier is trained on positive samples from OpenWebMath \citep{openwebmath}, while a heterogeneous array of unrelated webpages functions as negative samples. The classifier is then used to identify further positive examples within the CC, which are subjected to human curation for additional quality control. The enriched data subsequently refines the classifier’s training, further enhancing its effectiveness. Evaluation demonstrates that this large-scale dataset is robust: DeepSeekMath-Base 7B secures 64.2\% on GSM8K \citep{gsm8k} and 36.2\% on the MATH competition benchmark \citep{MATH}, surpassing Minerva 540B \citep{minerva}. Moreover, due to its multilingual composition, the DeepSeekMath Corpus yields improved results on Chinese mathematical assessment datasets \citep{wei2023cmath,agieval}. The methodological framework presented here offers a foundation for ongoing research in mathematical language processing, highlighting substantial opportunities for future advancements.

DeepSeekMath-Base is built upon DeepSeek-Coder-Base-v1.5 7B \citep{deepseek-coder}, as our findings suggest that initializing from a code-oriented model yields superior outcomes relative to beginning with a generic language model. Notably, training with mathematical data not only strengthens the model’s mathematical proficiency but also enhances its general reasoning abilities, as evidenced by improvements on the MMLU \citep{mmlu} and BBH \citep{bbh} benchmarks.

Subsequent to pre-training, we refine DeepSeekMath-Base through mathematical instruction tuning, leveraging datasets for chain-of-thought \citep{cot}, program-of-thought \citep{pot,pal}, and tool-integrated reasoning \citep{tora}. The resultant DeepSeekMath-Instruct 7B model surpasses all other models in the 7B class, achieving performance comparable to open-source instruction-tuned models at the 70B scale.

In addition, we present Group Relative Policy Optimization (GRPO), a new variant of the Proximal Policy Optimization (PPO) reinforcement learning (RL) framework \citep{schulman2017proximal}. Unlike standard PPO, GRPO omits the need for a critic model by inferring the baseline directly from group-level scores, thereby greatly reducing computational requirements. Applying GRPO to a selected subset of English instruction tuning data yields significant gains over the already robust DeepSeekMath-Instruct, with marked improvements across both in-domain (GSM8K: 82.9\% $\rightarrow$ 88.2\%, MATH: 46.8\% $\rightarrow$ 51.7\%) and out-of-domain tasks (for instance, CMATH: 84.6\% $\rightarrow$ 88.8\%) during reinforcement learning.

We also provide a unified framework that contextualizes various methods such as Rejection Sampling Fine-Tuning (RFT) \citep{yuan2023scaling}, Direct Preference Optimization (DPO) \citep{dpo}, PPO, and GRPO within a common perspective. Through this lens, all these techniques can be categorized as instances of either direct or streamlined RL approaches. To elucidate key elements of this paradigm, we conduct a comprehensive suite of experiments, examining, for example, online versus offline training, outcome versus process supervision, and single-turn versus iterative RL, among other aspects. Finally, we articulate the reasons behind the reinforcement learning improvements seen in instruction-tuned models and outline prospective avenues for advancing RL effectiveness within this unified schema.

\subsection{Contributions}
This work presents advancements in scalable math pre-training and a thorough investigation of reinforcement learning strategies.

\textbf{Math Pre-Training at Scale}

\begin{itemize}[topsep=0pt]
    \item We demonstrate that publicly available Common Crawl data encompasses a wealth of mathematical information. Utilizing a carefully engineered data selection pipeline, we curate the DeepSeekMath Corpus, a robust dataset containing 120B tokens sourced from web pages with substantial mathematical content. This dataset is nearly 7 times larger than the math-oriented web data employed by Minerva \citep{minerva} and approximately 9 times larger than the recently introduced OpenWebMath \citep{openwebmath}.

    \item Our base model, DeepSeekMath-Base 7B, delivers results comparable to those of Minerva 540B \citep{minerva}, suggesting that sheer parameter count is not the sole determinant of mathematical reasoning performance. Models of a modest scale, when pre-trained with high-quality data, can also attain strong results.

    \item We present insights from our math-specific training experiments. Pre-training on code before embarking on math training enhances a model’s capability to tackle mathematical tasks, regardless of whether tools are employed. This addresses the longstanding question: \textit{does code training improve reasoning abilities?} Our results suggest it does, especially concerning mathematical reasoning.

    \item Training models on arXiv publications—despite its popularity in math-focused works—fails to yield substantial gains across all mathematical benchmarks evaluated in this study.
\end{itemize}

\textbf{Exploration and Analysis of Reinforcement Learning}

\begin{itemize}[topsep=0pt]
    \item We present Group Relative Policy Optimization (GRPO), a reinforcement learning approach that prioritizes efficiency and effectiveness. By omitting the critic component and instead deriving baselines directly from group-level scores, GRPO drastically lowers computational demands relative to Proximal Policy Optimization (PPO).
    \item Our results show that applying GRPO notably improves the performance of our instruction-tuned system, DeepSeekMath-Instruct, utilizing only instruction-tuning datasets. In addition, we detect gains in generalization to out-of-domain problems throughout the reinforcement learning stage.
    \item We establish a cohesive framework that brings together various algorithms, such as RFT, DPO, PPO, and GRPO. To thoroughly examine the key attributes of this framework, we conduct diverse experimental comparisons, including online versus offline learning, outcome-oriented versus process-focused supervision, as well as contrasts between single-turn and iterative reinforcement learning.
    \item Leveraging this unified framework, we investigate the core factors underpinning the success of reinforcement learning and discuss several promising avenues for developing even more robust reinforcement learning methodologies for LLMs.
\end{itemize}

\subsection{Summary of Evaluations and Metrics}

\begin{itemize}[topsep=0pt]
    \item \textbf{English and Chinese Mathematical Reasoning}: Our evaluation strategy encompasses both English and Chinese datasets, spanning a range of mathematical tasks from elementary school up to university difficulty. The English datasets examined include GSM8K \citep{gsm8k}, MATH \citep{MATH}, SAT \citep{llemma}, OCW Courses \citep{minerva}, and MMLU-STEM \citep{mmlu}; the Chinese datasets assessed comprise MGSM-zh \citep{mgsm}, CMATH \citep{wei2023cmath}, Gaokao-MathCloze \citep{agieval}, and Gaokao-MathQA \citep{agieval}. We assess both the models' capabilities in producing self-contained textual solutions without external tools, and their proficiency in solving problems programmatically via Python.

    On the English tasks, DeepSeekMath-Base attains results comparable to those of the proprietary Minerva 540B \citep{minerva}, and exceeds the accuracy of all existing open-source base models—including Mistral 7B \citep{mistral} and Llemma-34B \citep{llemma}—regardless of the extent of prior math-centric pre-training, often by notable margins. On Chinese datasets, DeepSeekMath-Base excels even further, a performance likely attributed to our strategy of augmenting pre-training data with high-quality multilingual mathematical resources rather than restricting to English-only content as done in past work \citep{minerva,llemma}. Augmenting this with instruction-based tuning and reinforcement learning, the subsequent DeepSeekMath-Instruct and DeepSeekMath-RL set new records, crossing 50\% accuracy for the challenging MATH benchmark among all open-source systems for the first time.
\end{itemize}

\item \textbf{Formal Mathematics}: DeepSeekMath-Base is assessed on its ability to perform informal-to-formal theorem translation, utilizing the miniF2F dataset \citep{minif2f} in conjunction with the Isabelle proof assistant \citep{isabelle}, as described in \citep{dsp_proof}. The model exhibits strong autoformalization performance in few-shot scenarios.
\item \textbf{Natural Language Understanding, Reasoning, and Code}: To provide a thorough evaluation of general linguistic comprehension, reasoning ability, and code generation, we test DeepSeekMath-Base across several benchmarks: the Massive Multitask Language Understanding (MMLU) benchmark \citep{mmlu}, which spans 57 multiple-choice subject domains; BIG-Bench Hard (BBH) \citep{bbh}, composed of 23 tasks predominantly involving multi-step reasoning; as well as HumanEval \citep{codex} and MBPP \citep{mbpp}, standard assessments for programming language models. Results indicate that math pre-training enhances both language reasoning and understanding.

\section{Math Pre-Training}

\subsection{Data Collection and Decontamination}
Here, we detail the procedure for constructing the DeepSeekMath Corpus from Common Crawl. Figure \ref{fig:data_collect} illustrates our iterative framework for methodically expanding a large-scale mathematics corpus sourced from Common Crawl, starting with a small, curated seed dataset of mathematical texts. This pipeline can likewise be adapted for use in related domains such as code.

Our initial seed corpus is OpenWebMath \citep{openwebmath}, a collection of carefully selected mathematical texts from the web. Using these data, we train a fastText model \citep{joulin2016fasttext} to identify and retrieve additional mathematical webpages with similar characteristics. For training, we randomly sample 500,000 positive instances from OpenWebMath and select an equal number of negative examples from Common Crawl. Training is performed using an open-source fastText implementation\footnote{\url{https://fasttext.cc}}, configured with a vector dimension of 256, learning rate of 0.1, maximum word n-gram length of 3, a minimum of 3 occurrences per word, and 3 training epochs. To decrease the data volume from the full Common Crawl, we use deduplication techniques based on URLs and near-duplicate content, which reduces the corpus to 40B HTML webpages. The trained fastText model is then used to retrieve mathematical web content from the deduplicated collection. To ensure quality, retrieved pages are scored using fastText model predictions and only the highest-ranking examples are kept. To determine the appropriate dataset size, we conduct pre-training studies with the top 40B, 80B, 120B, and 160B tokens. In the initial iteration, we opt to retain the top 40B tokens.

Following the initial round of data collection, a significant number of mathematical web pages remain missing. This is primarily due to the limited diversity among positive examples used to train the fastText model. To address this, we expand the seed corpus by incorporating further mathematical web sources, thereby enabling enhanced optimization of the fastText model. Our methodology begins with segmenting the entire Common Crawl dataset into distinct domains, defined by web pages that share an identical base URL. For each domain, we determine the proportion of pages that were retrieved in the initial iteration. Any domain in which more than 10\% of pages have been collected is categorized as math-related (for instance, \url{mathoverflow.net}). 

Within these math-oriented domains, we manually annotate those URLs that host mathematical content (e.g., \url{mathoverflow.net/questions}). Web pages associated with these annotated URLs that were overlooked in the initial crawl are subsequently added to our seed set. Through this process, we expand the set of positive examples available for fastText model training, resulting in an improved model with higher recall for mathematical web pages in subsequent rounds. This iterative process is conducted four times, ultimately producing a collection of 35.5M math web pages, comprising a total of 120B tokens. By the fourth round, we observe that approximately 98\% of web pages had already been obtained during the third round, prompting us to conclude the data gathering process.

To prevent overlap with evaluation benchmarks, we adopt the filtering approach outlined in \cite{deepseek-coder}, ensuring that web pages containing content from English mathematical benchmarks like GSM8K \citep{gsm8k} and MATH \citep{MATH}, as well as Chinese benchmarks such as CMATH \citep{wei2023cmath} and AGIEval \citep{agieval}, are removed. Specifically, any textual segment that shares an exact 10-gram sequence with any part of an evaluation benchmark is excluded from our math training set. For benchmark samples that contain fewer than 10 grams but at least 3, we likewise apply strict string matching to remove potentially contaminated content.

\subsection{Validating the Quality of the DeepSeekMath~Corpus}

We conduct pre-training experiments to assess the quality of the DeepSeekMath~Corpus in comparison to other recently available mathematical corpora:

\begin{itemize}[topsep=0pt]
    \item \textbf{MathPile} \citep{mathpile}: This multi-source dataset contains 8.9B tokens drawn from a mix of sources, including textbooks, Wikipedia, ProofWiki, CommonCrawl, StackExchange, and arXiv—with arXiv accounting for more than 85\% of the content;
    \item \textbf{OpenWebMath} \citep{openwebmath}: Constructed by filtering CommonCrawl for mathematics-relevant web data, this corpus encompasses 13.6B tokens;
    \item \textbf{Proof-Pile-2} \citep{llemma}: This collection includes OpenWebMath, 10.3B tokens of mathematical code from AlgebraicStack, and 28.0B tokens from arXiv articles. In experiments involving Proof-Pile-2, we adopt the arXiv:Web:Code ratio of 2:4:1 as recommended by \cite{llemma}.
\end{itemize}

\subsubsection{Training Setting}
\label{sec:quality-policy}
We conduct mathematical training on a 1.3B-parameter general-purpose language model architecture, consistent with the DeepSeek LLMs \citep{deepseek-llm}, and refer to this as DeepSeek-LLM 1.3B. The model is independently fine-tuned on each mathematical dataset for 150B tokens. All training runs utilize the efficient and lightweight HAI-LLM framework \citep{haillm}. Adopting the procedures outlined for DeepSeek LLMs, optimization is performed using AdamW \citep{adamW} with hyperparameters $\beta_1=0.9$, $\beta_2=0.95$, and $\mathrm{weight\_decay}=0.1$. The learning rate schedule is multi-stage: after a 2,000-step warmup to maximum rate, it is reduced to 31.6\% of the peak after 80\% of total steps and then further down to 10.0\% of the peak rate by the final 10\% of training. The peak learning rate is set to 5.3e-4. We use a batch size of 4M tokens and a context window of 4K tokens.

\subsubsection{Evaluation Results}

\textbf{The DeepSeekMath~Corpus demonstrates exceptional quality, provides multilingual mathematical coverage, and represents the largest available dataset.}

\begin{itemize}[topsep=0pt]
    \item \textbf{High-quality}:
    We assess model performance across 8 mathematical benchmarks employing few-shot chain-of-thought prompting \cite{cot}. Table \ref{tab:corpora_comparison} indicates that the model trained on DeepSeekMath~Corpus surpasses others in downstream accuracy. Figure \ref{fig:corpora_comparisons} further illustrates that, at 50B tokens (corresponding to a full pass over Proof-Pile-2), the DeepSeekMath Corpus-trained model outperforms the Proof-Pile-2 model, suggesting that DeepSeekMath~Corpus offers higher average data quality.
    \item \textbf{Multilingual}:
    The DeepSeekMath~Corpus contains data across a range of languages, with English and Chinese as the most prominent. As detailed in Table \ref{tab:corpora_comparison}, leveraging DeepSeekMath~Corpus for training significantly advances mathematical reasoning in both English and Chinese. In contrast, existing corpora, which focus largely on English, show only modest gains or even adverse effects in Chinese mathematical reasoning.
    \item \textbf{Large-scale}:
    The DeepSeekMath~Corpus exceeds existing mathematical corpora in scale by a significant margin. According to Figure \ref{fig:corpora_comparisons}, DeepSeek-LLM 1.3B, when trained on DeepSeekMath~Corpus, experiences a more rapid ascent in performance and achieves sustained improvement. Meanwhile, alternative corpora, which are much smaller and cycled through multiple epochs during training, produce models whose accuracy quickly stagnates.
\end{itemize}

\subsection{Training and Evaluating DeepSeekMath-Base 7B}

This section presents DeepSeekMath-Base 7B, a foundational model designed with advanced mathematical reasoning capabilities. The model is initialized from DeepSeek-Coder-Base-v1.5 7B \citep{deepseek-coder} and undergoes training over 500B tokens. The composition of the training corpus is as follows: 56\% originates from the DeepSeekMath Corpus, 4\% is sourced from AlgebraicStack, 10\% is from arXiv, 20\% comprises Github code, while the final 10\% includes natural language material from Common Crawl in both English and Chinese. The training configuration aligns primarily with that described in Section \ref{sec:quality-policy}, with two modifications: the learning rate is capped at 4.2e-4 and a batch size of 10M tokens is employed.

To thoroughly evaluate DeepSeekMath-Base 7B's mathematical proficiency, we analyze its performance on generating complete mathematical solutions autonomously, leveraging external tools for problem-solving, and executing formal theorem proving. In addition to its mathematical evaluation, we also assess the model’s general capabilities, such as its competence in natural language understanding, logical reasoning, and coding.

\paragraph{Mathematical Problem Solving with Step-by-Step Reasoning}

We assess DeepSeekMath-Base's effectiveness in addressing mathematical challenges using few-shot chain-of-thought prompting \citep{cot}. Evaluations are carried out across eight distinct benchmarks in both English and Chinese. These tasks include quantitative reasoning assessments (such as GSM8K \citep{gsm8k}, MATH \citep{MATH}, and CMATH \citep{wei2023cmath}) and various multiple-choice tasks (such as MMLU-STEM \citep{mmlu} and Gaokao-MathQA \citep{agieval}), spanning a broad array of mathematical topics from basic to university-level intricacy.

As indicated in Table \ref{tab:base_math_cot}, among all open-source base models—including the popular Mistral 7B \citep{mistral} and the more recent, mathematically-trained Llemma 34B \citep{llemma}—DeepSeekMath-Base 7B achieves the highest scores across eight benchmark datasets. Remarkably, for the MATH dataset, which assesses competition-level mathematical ability, DeepSeekMath-Base exceeds the top open-source base models by a margin exceeding 10\% absolute, and even outperforms the substantially larger Minerva 540B \citep{minerva}—a proprietary model that is 77 times larger, derived from PaLM \citep{palm}, and further trained on mathematical data.

\paragraph{Mathematical Problem Solving with Tool Use} For program-aided reasoning in mathematics, we apply few-shot program-of-thought prompting \citep{pot,pal} on GSM8K and MATH. In this setup, models tackle problems by generating Python programs, making use of libraries like \textit{math} and \textit{sympy} to handle complex calculations. The final answer is determined by executing the produced code. Results in Table \ref{tab:base_math_pot_proof} demonstrate that DeepSeekMath-Base 7B outperforms the previously leading open-source model, Llemma 34B.

\paragraph{Formal Mathematics}

Automating the creation of formal proofs improves the correctness and efficiency of mathematical reasoning, drawing increasing interest. DeepSeekMath-Base 7B is evaluated on the informal-to-formal proving task defined in \citep{dsp_proof}, which requires converting an informal description, its formal translation, and an informal argument into a formal proof. We conduct experiments on miniF2F \citep{minif2f}, a standard benchmark for formalization of Olympiad-grade math, by prompting the model with several few-shot examples to construct formal Isabelle proofs for each item. Adopting the approach in \cite{dsp_proof}, the model generates a sketch of the proof, with Sledgehammer \citep{sledgehammer} subsequently used to automatically complete the remaining proof steps. Table \ref{tab:base_math_pot_proof} shows that DeepSeekMath-Base 7B delivers robust results for proof autoformalization.

\paragraph{Natural Language Understanding, Reasoning, and Code}

To assess natural language comprehension, we use MMLU \citep{mmlu}, for reasoning we employ BBH \citep{bbh}, and coding abilities are evaluated using HumanEval \citep{codex} and MBPP \citep{mbpp}. As summarized in Table \ref{tab:understanding_reasoning_code}, DeepSeekMath-Base 7B achieves considerable improvements on MMLU and BBH compared to its predecessor, DeepSeek-Coder-Base-v1.5 \citep{deepseek-coder}, indicating that mathematical training enhances the model’s language and reasoning capacities. By continuing pretraining with code data, DeepSeekMath-Base 7B also sustains the coding performance of DeepSeek-Coder-Base-v1.5 across HumanEval and MBPP. In summary, on all three reasoning and coding benchmarks, DeepSeekMath-Base 7B outperforms the general-purpose Mistral 7B \citep{mistral} by a significant margin.

\section{Supervised Fine-Tuning}

\subsection{SFT Data Curation}

We develop an instruction-tuning dataset for mathematics that encompasses both English and Chinese language tasks, drawing from a spectrum of mathematical disciplines and difficulty levels. Each problem is matched with its respective solution, annotated in chain-of-thought (CoT) \citep{cot}, program-of-thought (PoT) \citep{pot,pal}, or tool-integrated reasoning formats \citep{tora}. The complete dataset comprises 776,000 training instances.

\begin{itemize}[topsep=0pt]
    \item \textbf{English mathematical datasets}: We enrich GSM8K and MATH problems by annotating them with tool-integrated solutions, and include selected items from MathInstruct \citep{MathInstruct} and the Lila-OOD training data \citep{lila}, which provide problems solved using CoT and PoT methods. The English datasets span a wide variety of mathematical topics, including but not limited to algebra, probability, number theory, calculus, and geometry.
    \item \textbf{Chinese mathematical datasets}: Our collection contains K-12 mathematics problems in Chinese, distributed across 76 specific subfields such as linear equations. These are annotated with both chain-of-thought and tool-integrated solution approaches.
\end{itemize}

\subsection{Training and Evaluating DeepSeekMath-Instruct 7B}

Here, we detail the mathematical instruction-tuning procedure for DeepSeekMath-Instruct 7B, which builds upon DeepSeekMath-Base. Training data samples are concatenated at random up to a maximum input sequence of 4,000 tokens. The fine-tuning process is carried out for 500 iterations using a batch size of 256 and a fixed learning rate of 5e-5.

The mathematical competence of the resulting models is assessed with and without external tool assistance, using four quantitative reasoning benchmarks covering both English and Chinese languages. We conduct comparisons against the most advanced models available at the time, including:

\begin{itemize}[topsep=0pt]
    \item \textbf{Closed-source models}: This category comprises (1) the GPT series, particularly GPT-4 \citep{gpt4} and GPT-4 Code Interpreter~\footnote{\url{https://openai.com/blog/chatgpt-plugins\#\#code-interpreter}}, recognized for their state-of-the-art performance, (2) Gemini Ultra and Pro \citep{gemini}, (3) Inflection-2 \citep{inflection-2}, (4) Grok-1~\footnote{\url{https://x.ai/model-card}}, and additional prominent models recently introduced by Chinese technology companies including (5) Baichuan-3~\footnote{\url{https://www.baichuan-ai.com}}, and (6) the most current GLM-4 release~\footnote{\url{https://open.bigmodel.cn/dev/api\#glm-4}}
    
    from the GLM model family \citep{glm}.
\end{itemize}

Most of these models are designed for general applications and have typically been subject to various alignment processes.

\item \textbf{Open-source models} span both general-purpose and math-specialized systems. General-purpose models include (1) DeepSeek-LLM-Chat 67B \citep{deepseek-llm}, (2) Qwen 72B \citep{qwen}, (3) SeaLLM-v2 7B \citep{seallm}, and (4) ChatGLM3 6B \citep{chatglm3}. Models optimized for mathematical tasks feature (5) InternLM2-Math 20B~\footnote{\url{https://github.com/InternLM/InternLM-Math}}, an extension of InternLM2 that incorporates math-specific training with additional instruction tuning; (6) Math-Shepherd-Mistral 7B, which enhances Mistral 7B \citep{mistral} via PPO training \citep{schulman2017proximal} guided by a process-supervised reward function; (7) the WizardMath series \citep{wizardmath}, which advances mathematical reasoning within Mistral 7B and Llama-2 70B \citep{llama2} by utilizing evolve-instruct (an approach combining AI-generated instructions with instruction tuning) and PPO training on data primarily collected from GSM8K and MATH; (8) MetaMath 70B \citep{metamath}, representing Llama-2 70B after fine-tuning on an expanded GSM8K and MATH dataset; (9) ToRA 34B \cite{tora}, derived from CodeLlama 34B and fine-tuned for mathematical problem-solving with integrated tools; and (10) MAmmoTH 70B \citep{MathInstruct}, which involves instruction-tuning Llama-2 70B using MathInstruct.

As indicated in Table \ref{tab:sft_rl_math}, in the evaluation scenario prohibiting tool use, DeepSeekMath-Instruct 7B exhibits robust stepwise reasoning capabilities. Specifically, on the challenging, competition-grade MATH dataset, our model outperforms all open-source counterparts and surpasses most closed-source alternatives (such as Inflection-2 and Gemini Pro) by a margin of at least 9\% in absolute terms. This performance holds even relative to significantly larger models (such as Qwen 72B) and those that have undergone extensive math-oriented reinforcement learning, like WizardMath-v1.1 7B. DeepSeekMath-Instruct also demonstrates parity with the proprietary Chinese models GLM-4 and Baichuan-3 on the MATH benchmark, but it has yet to exceed the accuracy levels of GPT-4 and Gemini Ultra.

Allowing models to leverage both natural language reasoning and programmatic tool use for solutions, DeepSeekMath-Instruct 7B nearly attains 60\% accuracy on MATH, outperforming all known open-source models. For additional benchmarks, its performance closely matches that of DeepSeek-LLM-Chat 67B, a former state-of-the-art that has tenfold greater parameter count.

\section{Reinforcement Learning}

\subsection{Group Relative Policy Optimization} Reinforcement learning (RL) methods have consistently enhanced the mathematical reasoning performance of LLMs in post-Supervised Fine-Tuning (SFT) stages \citep{wang2023math,wizardmath}. Here, we detail our proposed RL technique, Group Relative Policy Optimization (GRPO), which emphasizes both efficiency and effectiveness.

\subsubsection{From PPO to GRPO} Proximal Policy Optimization (PPO) \citep{schulman2017proximal}, an actor-critic algorithm, has become standard for RL-based LLM fine-tuning \citep{ouyang2022training}. PPO focuses on maximizing the following clipped surrogate objective for model optimization:

\begin{equation}
\footnotesize
    \mathcal{J}_{PPO}(\theta) = \mathbb{E}{[q \sim P(Q), o \sim \pi_{\theta_{old}}(O|q)]} \frac{1}{|o|} \sum_{t=1}^{|o|} \min \left[ \frac{\pi_\theta(o_{t} | q, o_{<t})}{\pi_{\theta_{old}}(o_{t} | q, o_{<t})} A_{t}, \text{clip} \left( \frac{\pi_\theta(o_{t} | q, o_{<t})}{\pi_{\theta_{old}}(o_{t} | q, o_{<t})}, 1 - \varepsilon, 1 + \

In this context, $\pi_{\theta}$ and $\pi_{\theta_{old}}$ refer to the present and previous iterations of the policy models, respectively, while $q$ and $o$ denote the questions and corresponding outputs sampled from the dataset of questions as well as from the older policy $\pi_{\theta_{old}}$. The hyper-parameter $\varepsilon$ serves as a clipping threshold in PPO to help maintain stable training dynamics. The term $A_t$ stands for the advantage function, estimated using Generalized Advantage Estimation (GAE) \citep{gae}, which utilizes the sequence of rewards $\{r_{\ge t}\}$ alongside a learned value function $V_{\psi}$. Consequently, PPO requires joint training of a value function together with the policy network. To prevent the policy from excessively optimizing with respect to the reward model, the established practice is to incorporate a token-level KL-divergence penalty against a reference policy within the reward for every token \citep{ouyang2022training}, that is,

\begin{equation}
    r_{t} = r_\varphi(q, o_{\le t}) - \beta \log\frac{\pi_{\theta}(o_{t}|q, o_{<t})}{\pi_{ref}(o_{t}|q, o_{<t})},
\label{eq:PPO-reward}
\end{equation}

where $r_\varphi$ designates the reward function, $\pi_{ref}$ identifies the reference model (typically set as the initial SFT model), and $\beta$ determines the weight of the KL penalty.

Since the value function in PPO is generally another neural network of similar complexity to the policy model itself, it significantly increases resource usage, both in terms of memory and computation. Furthermore, the value function operates as a baseline within the advantage calculation for variance control during reinforcement learning. In the context of large language models (LLMs), however, it is common to allocate a reward only to the final token via the reward model, thereby complicating the process of accurately training a value function on a per-token basis. To resolve these issues, we introduce Group Relative Policy Optimization (GRPO), depicted in Figure \ref{fig:grpo}, which circumvents the need for training a separate value network as done in PPO. GRPO instead utilizes the mean reward obtained from several sampled outputs responding to a single question as the baseline. More precisely, given a question $q$, GRPO draws a group of outputs $\{o_1, o_2, \cdots, o_G\}$ from the earlier policy $\pi_{\theta_{old}}$, and proceeds to train the policy by maximizing this objective:

\begin{equation}
\footnotesize
\begin{split}
    \mathcal{J}_{GRPO}(\theta) &= \mathbb{E}{[q \sim P(Q), \{o_i\}_{i=1}^G \sim \pi_{\theta_{old}}(O|q)]}  \\
    & \frac{1}{G}\sum_{i=1}^G\frac{1}{|o_i|} \sum_{t=1}^{|o_i|} \left\{ \min \left[ \frac{\pi_\theta(o_{i,t} | q, o_{i,<t})}{\pi_{\theta_{old}}(o_{i,t} | q, o_{i,<t})} \hat{A}_{i,t}, \text{clip} \left( \frac{\pi_\theta(o_{i,t} | q, o_{i,<t})}{\pi_{\theta_{old}}(o_{i,t} | q, o_{i,<t})}, 1 - \varepsilon, 1 + \varepsilon \right)  \hat{A}_{i,t} \right] - \beta \mathbb{D}_{KL}\left[\pi_{\theta} || \pi_{ref}\right]\right\} ,
\end{split}
\label{eq:GRPO-obj}
\end{equation}

Here, $\varepsilon$ and $\beta$ act as hyper-parameters, while $\hat{A}_{i,t}$ denotes the advantage computed from only the relative rewards among outputs in the same group—a process discussed in more detail in later subsections. This group-based, relative computation of advantages in GRPO is naturally compatible with the inherently comparative design of most reward models, which are trained on datasets composed of output comparisons given identical prompts. Importantly, instead of integrating the KL penalty within the reward itself, GRPO imposes regularization by directly adding the KL divergence between the current policy and the reference policy to the loss, thus simplifying the computation of $\hat{A}_{i,t}$. Additionally, rather

\begin{algorithm}[t]
  \small
  \caption{Iterative Group Relative Policy Optimization}
  \textbf{Input} initial policy model $\pi_{\theta_{\text{init}}}$; reward models $r_\varphi$; task prompts $\mathcal{D}$; 
  hyperparameters $\varepsilon$, $\beta$, $\mu$
  \begin{algorithmic}[1]
    \State Initialize policy model $\pi_\theta \leftarrow \pi_{\theta_{\text{init}}}$
    \For{iteration = 1, \dots, I}
       \State Set reference model $\pi_{ref} \leftarrow \pi_{\theta}$
      \For{step = 1, \dots, M}
      \State Draw a mini-batch $\mathcal{D}_b$ from $\mathcal{D}$
      \State Save the previous policy parameters $\pi_{\theta_{old}} \leftarrow \pi_{\theta}$ 
      \State For each question $q \in \mathcal{D}_b$, generate $G$ outputs $\{o_i\}_{i=1}^G \sim \pi_{\theta_{old}} (\cdot \mid q) $
      \State Compute the reward values $\{r_i\}_{i=1}^{G}$ for each $o_i$ via $r_{\varphi}$
      \State Calculate $\hat{A}_{i,t}$ for the $t$-th token in $o_i$ using the group relative advantage estimator
      \For{GRPO iteration = 1, \dots, $\mu$}
        \State Update $\pi_{\theta}$ by maximizing the GRPO objective (Equation \ref{eq:GC-GRPO})
      \EndFor
    \EndFor 
    \State Refine $r_\varphi$ through ongoing training with a replay buffer. 
    \EndFor 
  \end{algorithmic}
  \textbf{Output} $\pi_\theta$
  \label{alg:iter-grpo}
\end{algorithm}

\subsubsection{Outcome Supervision RL with GRPO} Precisely, for a given question $q$, the old policy model $\pi_{\theta_{old}}$ samples a set of outputs $\{o_1, o_2, \cdots, o_G\}$. Each output is assigned a reward through the reward model, resulting in a reward vector $\mathbf{r} = \{r_1, r_2, \cdots, r_G\}$. These rewards are standardized by subtracting the group mean and dividing by the group standard deviation. Outcome supervision employs this normalized reward at the completion of every $o_i$ and attributes this same normalized value as the advantage $\hat{A}_{i, t}$ to every token in $o_i$, formally, $\hat{A}_{i, t} = \widetilde{r}_i = \frac{r_i- {\rm mean}(\mathbf{r})}{{\rm std}(\mathbf{r})}$. The policy is then trained by maximizing the criterion in equation (\ref{eq:GRPO-obj}).

\subsubsection{Process Supervision RL with GRPO} Since outcome supervision restricts feedback to the end of the sequence, it can be less effective for intricate mathematical reasoning. To address this, following the approach of \cite{wang2023math}, we also consider process supervision, which distributes rewards across each reasoning stage. For a question $q$ and $G$ generated outputs $\{o_1, o_2, \cdots, o_G\}$, the process reward model evaluates every individual step, producing structured rewards: $\mathbf{R} = \{ \{r_1^{index(1)},\cdots,r_1^{index(K_1)}\}, \cdots,  \{r_G^{index(1)},\cdots,r_G^{index(K_G)}\} \}$, where $index(j)$ identifies the end token position of the $j$-th reasoning step, and $K_i$ is the number of steps within the $i$-th output. Normalization is again applied, yielding $\widetilde{r}_i^{index(j)} = \frac{r_i^{index(j)} - {\rm mean(\mathbf{R})}}{{\rm std(\mathbf{R})}}$.
With process supervision, the advantage assigned to each token is determined by summing all subsequent normalized step rewards, that is, $\hat{A}_{i, t} = \sum_{index(j) \ge t} \widetilde{r}_i^{index(j)}$,
and the policy is subsequently optimized by maximizing the objective outlined in equation (\ref{eq

\subsubsection{Iterative RL with GRPO} During the reinforcement learning process, the existing reward model might no longer provide adequate guidance to the evolving policy model. To address this, we investigate iterative RL utilizing GRPO. As illustrated in Algorithm \ref{alg:iter-grpo}, the iterative GRPO approach involves generating updated training datasets for the reward model by sampling outputs from the current policy model. The reward model is then further trained using a replay buffer that retains 10\% of previously encountered data. Subsequently, the policy model is used as the new reference model, and it is continually optimized using feedback from the updated reward model.

\subsection{Training and Evaluating DeepSeekMath-RL}

We apply reinforcement learning starting from the DeepSeekMath-Instruct 7B model. The RL training dataset includes chain-of-thought questions related to GSM8K and MATH extracted from the SFT dataset, comprising approximately 144K questions. In order to analyze the specific effects of reinforcement learning on benchmarks without any overlapping data during the RL stage, other SFT questions are omitted. The methodology for creating the reward model's training dataset follows \citep{wang2023math}. The initial reward model is trained on DeepSeekMath-Base 7B, using a learning rate of 2e-5. For GRPO, we assign the policy model a learning rate of 1e-6, set the KL divergence coefficient to 0.04, and generate $64$ samples per question. The maximum sequence length is limited to 1024 tokens, with a batch size also set to 1024. The policy model undergoes a single update after each exploration round. DeepSeekMath-RL 7B is evaluated on a suite of benchmarks in the same manner as DeepSeekMath-Instruct 7B. Here, GSM8K and MATH tasks featuring chain-of-thought reasoning are considered in-domain, while all remaining benchmarks are categorized as out-of-domain.

Table \ref{tab:sft_rl_math} presents a comparative analysis of both open-source and closed-source models, evaluated using chain-of-thought and tool-based reasoning across English and Chinese benchmarks. Our observations reveal the following: 1) DeepSeekMath-RL 7B achieves 88.2\% accuracy on GSM8K and 51.7\% on MATH when employing chain-of-thought reasoning. These results exceed those of all open-source models within the 7B to 70B parameter scale and outperform most closed-source counterparts. 2) Significantly, DeepSeekMath-RL 7B is trained solely on instruction-tuning data in the chain-of-thought format derived from GSM8K and MATH, with initialization from DeepSeekMath-Instruct 7B. Despite this limited training data, it demonstrates consistent performance gains over DeepSeekMath-Instruct 7B across all evaluation criteria, underscoring the strengths of reinforcement learning.

\section{Discussion}
In this section, we discuss insights gleaned from both our pre-training and reinforcement learning studies.

\subsection{Insights from Pre-Training}

Here, we recount our pre-training experience. Unless explicitly mentioned otherwise, the training protocols conform to the specifications in Section \ref{sec:quality-policy}. For clarity, any reference to the DeepSeekMath Corpus pertains to the 89B-token dataset generated during the second phase of our data collection.

\subsubsection{The Impact of Code Training on Mathematical Reasoning}

An often-cited but not conclusively established hypothesis posits that code training bolsters reasoning abilities. We contribute partial evidence to this claim in the context of mathematics: training on code data enhances mathematical reasoning proficiency in models, with or without tool assistance.

To assess the impact of code training on mathematical problem-solving, we conducted investigations under two distinct regimes: two-stage training and one-stage training.

\textbf{Two-Stage Training}

\begin{itemize}[topsep=0pt]
    \item \textbf{Code Training for 400B Tokens $\rightarrow$ Math Training for 150B Tokens}: DeepSeek-LLM 1.3B undergoes pretraining on 400B tokens of code, followed by a subsequent phase involving 150B tokens of math-specific data;
    \item \textbf{General Training for 400B Tokens $\rightarrow$ Math Training for 150B Tokens}: To provide a baseline, we alternatively pretrain on 400B general-domain tokens (sampled from DeepSeek-AI's large general corpus) instead of code in the first stage, allowing us to examine whether code token pretraining confers advantages over general tokens for mathematical reasoning skills.
\end{itemize}

\textbf{One-Stage Training}

\begin{itemize}[topsep=0pt]
    \item \textbf{Math Training for 150B Tokens}: The DeepSeek-LLM 1.3B model is trained exclusively on 150B math tokens;
    \item \textbf{Training on a mixture of 400B Code Tokens and 150B Math Tokens}: Recognizing that mathematical finetuning after code pretraining can harm coding capability, we further assess whether integrating code and math tokens within a single-phase training setup could sustain or enhance both mathematical reasoning and coding abilities, potentially counteracting catastrophic forgetting effects.
\end{itemize}

\paragraph{Results}
Tables \ref{tab:code-math} and \ref{tab:code-math-general-eval} summarize task performance under varying training schemes.

Engagement with code data prior to or alongside math tokens improves program-aided mathematical reasoning, evident across both two-stage and single-stage paradigms. Specifically, Table \ref{tab:code-math} indicates that initial code-based pretraining under the two-phase regime greatly strengthens the model’s proficiency at solving GSM8K and MATH tasks with Python. Additional math-oriented training in the latter phase amplifies these benefits. Furthermore, in one-stage approaches, simultaneous exposure to both code and math tokens notably diminishes the impact of catastrophic forgetting typically seen in sequential training, while still fostering gains in both coding (Table \ref{tab:code-math-general-eval}) and math reasoning with programmatic support (Table \ref{tab:code-math}).

Beyond tool-augmented contexts, pretraining with code also yields improvements in unaided mathematical reasoning: the model experiences clear early gains from code data in a two-stage trajectory, which subsequently expedites math learning for peak results. However, mixing code and math tokens together from the outset within a one-stage framework appears to impair the model’s capability to reason mathematically without code execution. This may be attributed to capacity constraints of the 1.3B DeepSeek-LLM, limiting its ability to simultaneously learn from both modalities at this smaller scale.

\subsubsection{ArXiv Papers Appear to Have Limited Impact on Enhancing Mathematical Reasoning}

Despite their prevalence as a source in mathematics-focused pre-training datasets \citep{minerva,gpt-f,llemma,mathpile}, there is a lack of in-depth evaluation regarding the efficacy of arXiv papers in bolstering mathematical reasoning abilities. Somewhat surprisingly, our experimental findings indicate that arXiv papers do not substantially benefit mathematical reasoning. We performed comparative studies on models of various capacities, specifically DeepSeek-LLM 1.3B and DeepSeek-Coder-Base-v1.5 7B \citep{deepseek-coder}, leveraging arXiv-based corpora that underwent distinct processing procedures:

\begin{itemize}[topsep=0pt]
    \item \textbf{MathPile} \citep{mathpile}: a dataset comprising 8.9 billion tokens, curated through cleaning and heuristic filtering processes, with scientific arXiv papers making up more than 85\% of its contents;
    \item \textbf{ArXiv-RedPajama} \citep{redpajama}: a full collection of arXiv LaTeX source files, where extraneous content such as preambles, comments, macros, and bibliographies are stripped, totaling 28.0 billion tokens.
\end{itemize}

For our investigation, DeepSeek-LLM 1.3B was trained for 150 billion tokens and DeepSeek-Coder-Base-v1.5 7B for 40 billion tokens independently on each arXiv dataset. Across the board, neither model demonstrated marked progress—or even showed regression—in mathematical reasoning skills when trained exclusively on arXiv corpora. Performance remained stagnant or declined on a range of mathematical tasks and benchmarks of varying complexity considered in our work, such as quantitative datasets like GSM8K and MATH (Table \ref{tab:arxiv-cot}), multiple-choice assessments such as MMLU-STEM (Table \ref{tab:arxiv-cot}), and formal mathematics evaluations like miniF2F (Table \ref{tab:arxiv_atp}).

Nevertheless, this interpretation is preliminary and subject to important caveats. Our study does not yet address the following open questions:

\begin{itemize}[topsep=0pt]
    \item How arXiv-based pre-training affects specialized mathematical activities not assessed here, including the transformation of formal proofs into informal statements;
    \item The interplay of arXiv data with other sources of mathematical data in a mixed dataset setting;
    \item Whether potential benefits of arXiv-sourced materials become pronounced at greater model scales.
\end{itemize}

Consequently, comprehensive investigation into these dimensions remains a direction for future research.

\subsection{Perspectives on Reinforcement Learning}
\subsubsection{Toward a Unified Framework} Here, we introduce a general framework for examining a variety of training algorithms, such as SFT, RFT, DPO, PPO, and GRPO. We also present empirical analysis to investigate key components of this unified perspective. In broad terms, the gradient with respect to the model parameters $\theta$ under a given training strategy can be formalized as:

\begin{equation}
    \nabla_{\theta}\mathcal{J}_{\textcolor{red}{\mathcal{A}}}(\theta) = \mathbb{E}[\underbrace{(q,o) \sim \textcolor{red}{\mathcal{D}}}_{Data \ Source}]\left( \frac{1}{|o|} \sum_{t=1}^{|o|}  \underbrace{GC_{{\mathcal{A}}}(q, o, t, \textcolor{red}{\pi_{{rf}}})}_{Gradient \ Coefficient}  \nabla_{\theta}\log \pi_{\theta}(o_t | q, o_{<t})\right).
\label{eq:objective}
\end{equation}

This formulation is characterized by three principal elements: 1) the \textit{Data Source $\mathcal{D}$}, which provides the basis for constructing the training set; 2) the \textit{Reward Function $\pi_{{rf}}$}, serving as the generator of feedback signals during training; and 3) the \textit{Algorithm $\mathcal{A}$}, which synthesizes information from the data source and reward function to compute the gradient coefficient $GC$, thereby modulating the strength of updates (penalization or reinforcement) applied to training samples. Based on this integrative framework, we compare and discuss several representative strategies:

\begin{itemize}
    \item  \textbf{Supervised Fine-tuning (SFT)}: This approach adapts a pretrained model by further training on manually curated SFT datasets.
    \item \textbf{Rejection Sampling Fine-tuning (RFT)}: RFT enhances the SFT-trained model by fine-tuning it on those responses, sampled from itself, that pass correctness filtering based on SFT queries.
    \item \textbf{Direct Preference Optimization (DPO)}: DPO optimizes beyond SFT by employing a pairwise preference-based objective, using additional sampled outputs to further refine the model.
    \item \textbf{Online Rejection Sampling Fine-tuning (Online RFT)}: Unlike standard RFT, Online RFT begins with the SFT-initialized policy and iteratively updates it via fine-tuning on newly generated outputs sampled online from the latest policy version.
    \item \textbf{PPO/GRPO}: In these methods, the policy is initialized using SFT and then strengthened through reinforcement, leveraging real-time samples generated from the evolving policy.
\end{itemize}

The elements comprising these methods are detailed in Table \ref{tab:data_policy}. For an expanded and in-depth explanation of the derivation process, see Appendix \ref{app:analysis-rl}.

\paragraph{Observation about Data Source} We categorize the sources of data as either online or offline sampling. Online sampling refers to gathering training data via the real-time exploration of the current policy model, whereas offline sampling obtains training data from outputs generated by the original SFT model. The RFT and DPO approaches are representative of the offline data paradigm, while Online RFT and GRPO rely on online data sampling.

Analysis presented in Figure \ref{fig:rl-analysis} reveals that Online RFT delivers a marked improvement over RFT across two evaluation benchmarks. In the initial phases of training, Online RFT exhibits similar performance to RFT; however, as training advances, Online RFT obtains a clear and decisive lead, highlighting the advantages associated with online data collection. This outcome is predictable: at the onset of training, there is only slight divergence between the actor and the SFT model, resulting in minimally varied samples. With continued training, distinctions in actor-sampled data become more pronounced, allowing real-time sampling to provide increasing benefits.

\paragraph{Observation about Gradient Coefficient} In these algorithms, the reward signal is transformed into a gradient coefficient for updating the model's parameters. Within our experimental framework, the reward function is categorized as either `Rule'—where assessment relies on answer correctness—or `Model'—in which a learned reward model scores responses, trained using data labeled by the rule-based approach. As illustrated by Equations \ref{eq:GC-RFT} and \ref{eq:GC-GRPO}, a pivotal distinction between GRPO and Online RFT emerges: GRPO modifies its gradient coefficient according to the reward assigned by the model, allowing the algorithm to intensify rewards or penalties in proportion to response quality. Online RFT, in contrast, does not implement this differentiation; it avoids penalization of incorrect responses and provides equivalent positive reinforcement to all correct answers, regardless of reward magnitude.

As illustrated in Figure \ref{fig:rl-analysis}, GRPO achieves better results than online RFT, underscoring the effectiveness of adjusting the coefficients for positive and negative gradients. Moreover, the performance of GRPO+PS outstrips that of GRPO+OS, which emphasizes the advantage of utilizing step-sensitive, fine-grained gradient coefficients. We further investigate the impact of iterative RL, conducting two iterations in our experiments. Figure \ref{fig:iter-rl} reveals that iterative RL yields notable performance gains, particularly after the initial round of iteration.

\subsubsection{Why RL Works?} In this study, reinforcement learning is performed using a selected portion of instruction tuning data, resulting in substantial improvements relative to the baseline instruction tuning model. To deepen the understanding of RL’s effectiveness, we measure both Pass@K and Maj@K accuracies for the Instruct and RL models across two evaluation benchmarks. As depicted in Figure \ref{fig:maj-pass}, reinforcement learning enhances Maj@K, while Pass@K remains largely unchanged. These results imply that RL primarily bolsters the robustness of the output distribution—\textbf{suggesting the observed gains stem from an increased likelihood of a correct response among the TopK predictions, rather than from intrinsic skill improvements.} A related observation is made by \citep{wang2023making}, who describe a \textbf{misalignment problem} within SFT models for reasoning tasks. Their work, as well as subsequent research \citep{yuan2023rrhf,song2023preference,wang2023making}, demonstrates that SFT models’ reasoning capabilities can be advanced through various preference alignment techniques.

\subsubsection{How to Achieve More Effective RL?}
\label{sec:effec_rl}
Our findings show that RL yields considerable success in mathematical reasoning settings. We also introduce a unified perspective, encompassing various representative training approaches under a common framework, where all techniques are classified as forms of either full or simplified RL. As formalized in Equation \ref{eq:objective}, this paradigm consists of three essential elements: Data Source, Algorithm, and Reward Function. We outline potential research opportunities for each component.

\paragraph{Data Source} The data source underpins all training approaches. In the RL framework, this refers to unlabeled problems for which the policy model generates responses. In our present work, only questions from instruction tuning, together with basic nucleus sampling, are used for generating outputs. We posit that this limitation may contribute to the observed improvements being restricted to Maj@K. Looking ahead, we plan to assess the RL pipeline using out-of-distribution prompts and \textbf{advanced sampling (decoding) strategies} such as tree-search-based methods \citep{yao2023tree}. Furthermore, \textbf{efficient inference techniques} \citep{xia-etal-2023-speculative,leviathan2023fast,kwon2023efficient,xia2024unlocking}, which are key to enhancing the policy model’s exploration capabilities, are anticipated to play a pivotal role as well.

\paragraph{Algorithms} Algorithms leverage both the data and the reward signals to compute gradient coefficients, which in turn adjust the model parameters. As derived from Equation \ref{eq:objective}, contemporary approaches almost universally \textbf{TRUST} the reward function’s feedback when modulating the conditional probability of specific tokens. Nevertheless, guaranteeing the constant reliability of this reward feedback is unattainable, particularly for tasks marked by high complexity. A case in point is the PRM800K datasets \citep{lightman2023let}; despite rigorous annotation by skilled annotators, around 20\% of the labels remain erroneous\footnote{\url{https://github.com/openai/prm800k/issues/12\#issuecomment-1728491852}}. In light of this, our focus turns to reinforcement learning algorithms capable of maintaining performance despite noisy reward inputs. We anticipate that such \textbf{WEAK-TO-STRONG} \citep{burns2023weak} alignment paradigms could drive transformative progress in the core mechanics of learning algorithms.

\paragraph{Reward Function} The reward function serves as the foundational training signal within RL, typically instantiated by a neural reward model. We identify three pivotal avenues for the development of reward models: 1) \textbf{Strengthening the generalization capabilities of reward models.} To meaningfully improve the behavior of LLMs, reward models must transfer effectively to out-of-distribution queries and advanced decoding scenarios; otherwise, reinforcement learning might only maintain the status quo of the model’s output distribution; 2) \textbf{Incorporating uncertainty estimates in reward modeling.} Capturing uncertainty can help bridge the gap between less reliable reward signals and algorithms designed for weak-to-strong improvement; 3) \textbf{Constructing high-quality process reward models efficiently.} These models are essential for supplying nuanced, fine-grained feedback during complex reasoning tasks \citep{lightman2023let,wang2023math}.

\section{Conclusion, Limitation, and Future Work}

This work introduces DeepSeekMath, a model that surpasses all publicly available models on the competition-level MATH benchmark and nears the efficacy of proprietary counterparts. DeepSeekMath~is based on the DeepSeek-Coder-v1.5 7B architecture and has undergone continual pre-training on 500B tokens, with 120B of those tokens consisting of mathematics-related data collected primarily from Common Crawl. Our comprehensive ablation studies indicate that web-derived content can serve as a rich source for high-quality mathematical datasets, whereas resources such as arXiv may be less advantageous than initially anticipated. We present Group Relative Policy Optimization (GRPO), a modification of Proximal Policy Optimization (PPO), which demonstrably enhances mathematical reasoning proficiency while consuming less memory. The results confirm GRPO's efficacy, showing marked improvements even when DeepSeekMath-Instruct 7B has already attained high benchmark scores. Furthermore, we offer a unified conceptual framework to situate these methods and highlight multiple promising directions for advancing reinforcement learning strategies.

Despite the notable advancements achieved by DeepSeekMath~in quantitative reasoning tasks, its performance remains limited in domains such as geometry and theorem proving, particularly when compared to leading closed-source models. For example, during preliminary tests, the model failed to solve problems involving triangles and ellipses, suggesting a possible bias in data selection during pre-training and fine-tuning stages. Moreover, owing to its smaller model size, DeepSeekMath~underperforms relative to GPT-4 in few-shot learning scenarios: while GPT-4 benefits from few-shot prompting, DeepSeekMath~demonstrates negligible performance variation between zero-shot and few-shot conditions. Looking ahead, we aim to enhance our data selection and curation pipeline for pre-training, targeting the assembly of higher-quality datasets. Additionally, we intend to pursue the directions proposed in Section \ref{sec:effec_rl} to realize more effective reinforcement learning methodologies for LLMs.

\end{document}